\documentclass{article}

\usepackage{amsmath,amssymb}
\usepackage{xurl}
\usepackage[hidelinks]{hyperref}
\setlength{\emergencystretch}{2em}

\input{command}

\title{Zero correlation length uses the physical-trace transfer in
Cirac--Perez-Garcia--Schuch--Verstraete 2017}
\author{}
\date{2026-06-08}

\begin{document}
\maketitle
\gapnote{local-correction}{resolved}

\section{Source statement}

For mixed states, Definition~4.2 of
Ref.~\cite{Cirac2016MPDO_arXiv}, with local label
\texttt{DefinitionZCL} in
\path{Papers/1606.00608/MPDO-22-12-17-2.tex:735--739}, defines a tensor $M$
generating an MPDO to have \emph{zero correlation length} (ZCL) through the
graphical equation in \path{MPDO_ZCL.png}. Let $M^{ij}$ denote the bond matrix
with ket index $i$ and bra index $j$. The diagram closes the ket and bra
physical legs of one MPDO tensor, so it represents the physical-trace transfer
\begin{align}
    \mathcal T_M
    &:=\sum_i M^{ii},
    \label{eq:zcl_physical_trace_transfer}
\end{align}
not the doubled-index completely positive map on bond matrices. The source
calls this condition ``the natural extension of the pure state case of
Theorem~\texttt{TheoremZCLPure}'' and states that correlation functions
computed from an MPDO with ZCL are length independent
\cite[lines~735--739]{Cirac2016MPDO_arXiv}.

The proof of Theorem~\texttt{TheoremZCLPure} in
Ref.~\cite{Cirac2016MPDO_arXiv}, at
\path{Papers/1606.00608/MPDO-22-12-17-2.tex:1248}, states that pure-state ZCL
is equivalent to
\begin{align}
    \mathbb E^2
    &=\mathbb E
    \label{eq:zcl_pure_transfer_idempotence}
\end{align}
for a tensor $A$ in canonical form. The renormalization step is
\begin{align}
    \mathcal E
    &\longmapsto\mathcal E^2,
    \label{eq:zcl_pure_renormalization_step}
\end{align}
as stated at
\path{Papers/1606.00608/MPDO-22-12-17-2.tex:1209}. A fixed point therefore
satisfies
\begin{align}
    \mathcal E
    &=\mathcal E^2.
    \label{eq:zcl_pure_fixed_point}
\end{align}
Canonical form normalizes the leading eigenvalue of the transfer operator to
one. Equation~(\ref{eq:zcl_pure_transfer_idempotence}) then states that
$\mathbb E$ projects onto its peripheral eigenvalue-one subspace.

\section{The gap}

Define the doubled-index completely positive transfer map by
\begin{align}
    \mathcal E_M(X)
    &=\sum_{i,j}M^{ij}X(M^{ij})^\ast.
    \label{eq:zcl_doubled_transfer_definition}
\end{align}
The older Lean definition \leanid{MPOTensor.IsZCL} instead requires
\begin{align}
    \mathcal E_M\circ\mathcal E_M
    &=\mathcal E_M.
    \label{eq:zcl_doubled_transfer_idempotence}
\end{align}
This is literal idempotence of the doubled-index completely positive transfer
map, not the transfer drawn in Definition~4.2 of
Ref.~\cite{Cirac2016MPDO_arXiv}. The source first contracts the two physical
legs and then squares $\mathcal T_M$ from
Eq.~(\ref{eq:zcl_physical_trace_transfer}). Thus the older predicate is not
merely an unnormalized version of the source condition; it concerns a
different transfer object. Even up to scalar normalization, idempotence of
$\mathcal E_M$ classifies tensors through the doubled-index MPS transfer rather
than through the MPDO ZCL diagram.

\paragraph{Classification: resolved definition drift.}
The older formal predicate records doubled-index CP-transfer idempotence,
whereas the source definition uses the physical-trace transfer. The source
definition is now formalized by
\leanid{MPOTensor.IsPhysicalTraceIdempotent}, whose equation characterization
is
\begin{align}
    \mathcal T_M^2
    &=\mathcal T_M.
    \label{eq:zcl_literal_physical_idempotence}
\end{align}
The predicate \leanid{MPOTensor.IsSourceZCL} records a separate,
scale-invariant up-to-scalar relation. It is distinct from both the paper's
literal diagram and the doubled-index predicate.

The development witnesses this deviation through
\leanid{MPOTensor.exists_isLocalPurificationRFP_not_isZCL}. Take
$d=d_K=2$, $D'=D=1$, and
\begin{align}
    A
    &=\left[\tfrac{1}{\sqrt2},0,0,\tfrac{1}{\sqrt2}\right].
    \label{eq:zcl_purification_tensor}
\end{align}
This is an RFP MPS tensor with
\begin{align}
    E_A
    &=\operatorname{id}.
    \label{eq:zcl_purification_transfer_identity}
\end{align}
The induced MPDO tensor has
\begin{align}
    M^{00}
    &=M^{11}
    =\tfrac12,
    \label{eq:zcl_mixed_diagonal_entries}\\
    M^{01}
    &=M^{10}
    =0.
    \label{eq:zcl_mixed_offdiagonal_entries}
\end{align}
Its length-$N$ periodic operator is
\begin{align}
    \rho^{(N)}(M)
    &=\left(\tfrac12\right)^N I.
    \label{eq:zcl_maximally_mixed_family}
\end{align}
This normalized maximally mixed state has length-independent correlations and
therefore satisfies ZCL in the source's sense. Its physical-trace transfer is
\begin{align}
    \mathcal T_M
    &=M^{00}+M^{11}
    \label{eq:zcl_maximally_mixed_trace_sum}\\
    &=1,
    \label{eq:zcl_maximally_mixed_trace_identity}
\end{align}
so the graphical idempotence holds. By contrast, its doubled-index CP transfer
satisfies
\begin{align}
    \mathcal E_M
    &=\tfrac12\operatorname{id},
    \label{eq:zcl_maximally_mixed_doubled_transfer}\\
    \mathcal E_M\circ\mathcal E_M
    &=\tfrac14\operatorname{id}
    \label{eq:zcl_maximally_mixed_doubled_square}\\
    &\ne\mathcal E_M.
    \label{eq:zcl_maximally_mixed_doubled_failure}
\end{align}
Thus literal \leanid{MPOTensor.IsZCL} fails. The witness exposes the wrong
transfer object, not merely a missing scalar normalization.

\section{Consequence and elimination plan}

The source implication
$\mathrm{PRFP}\Rightarrow\mathrm{ZCL}$ at
\path{Papers/1606.00608/MPDO-22-12-17-2.tex:777} in
Ref.~\cite{Cirac2016MPDO_arXiv} cannot be represented by
\leanid{MPOTensor.IsZCL}: the source implication controls $\mathcal T_M$,
whereas that predicate controls $\mathcal E_M$. After this definition drift is
repaired, the printed implication to literal physical-trace idempotence is
itself refuted by
\leanid{MPOTensor.exists_isPRFP_isMPDO_physTraceTransfer_ne_zero_not_isPhysicalTraceIdempotent}.
Its witness satisfies global PRFP and MPDO positivity with
\begin{align}
    \mathcal T_M
    &\ne0,
    \label{eq:zcl_prfp_witness_nonzero_transfer}\\
    \mathcal T_M^2
    &\ne\mathcal T_M.
    \label{eq:zcl_prfp_witness_not_idempotent}
\end{align}
The stronger failure of the scale-invariant relation remains a separate
project result.

Two normalization conventions remain useful for corrected restricted
theorems:
\begin{enumerate}
    \item The source ZCL transfer uses
    $\mathcal T_M=\sum_iM^{ii}$ and allows scalar normalization at the level of
    the generated state family. Its nondegenerate up-to-scalar condition is
    \begin{align}
        \mathcal T_M
        &\ne0,
        \label{eq:zcl_source_transfer_nonzero}\\
        \mathcal T_M^2
        &=\lambda\mathcal T_M
        &&\text{for some $\lambda>0$}.
        \label{eq:zcl_source_transfer_quasi_idempotence}
    \end{align}
    The physical-trace-normalized representative is the case $\lambda=1$.
    \item The canonical physical-trace hypothesis retains exact idempotence
    in Eq.~(\ref{eq:zcl_literal_physical_idempotence}) only for source-labelled
    theorems whose hypotheses include the required canonical normalization.
    The doubled-index map $\mathcal E_M$ remains a separate completely positive
    transfer used in the formal fusion and algebra statements, not the
    source's ZCL transfer.
\end{enumerate}

The source definition and the scale-invariant relation are now both
formalized. The declaration
\leanid{MPOTensor.IsPhysicalTraceIdempotent} records the literal fixed-tensor
diagram, \leanid{MPOTensor.IsSourceZCL} records the nonzero up-to-scalar
relation, and \leanid{MPOTensor.IsZCL} retains the doubled-index CP-transfer
condition. These predicates must remain distinct.

For the renormalization condition in Definition~4.1 of
Ref.~\cite{Cirac2016MPDO_arXiv}, Proposition~\texttt{propsimple} is now proved
at the explicit normalization boundary: every positive-length ring is
positive semidefinite and has nonzero trace. The primary source-facing
declarations are \leanid{MPOTensor.IsPhysicalTraceIdempotent} for
Definition~4.2 and
\leanid{MPOTensor.isPhysicalTraceIdempotent_of_isRFPViaTS} for its RFP
implication, together with
\leanid{MPOTensor.isSAL_of_isRFPViaTS_of_trace_ne_zero} for SAL. The
conjunction
\leanid{MPOTensor.isSourceZCL_and_isSAL_of_isRFPViaTS_of_trace_ne_zero}
has the broader scale-invariant conclusion and is a convenience theorem, not
the paper-facing statement.

The horizontal source-facing specialization uses
\leanid{MPOTensor.physTraceTransfer_sq_of_isRFPViaTS} together with
\leanid{MPOTensor.isSAL_of_isRFPViaTS}. The corresponding
\leanid{MPOTensor.IsSourceZCL} conjunction is again only a convenience result.
Positivity alone admits the zero tensor, so the trace condition cannot be
omitted. The printed global PRFP equivalence is refuted; the local-purification
equivalences remain valid only under their stated tensor-level restriction.

In Theorem~4.9 of Ref.~\cite{Cirac2016MPDO_arXiv}, the literal implication
(iv)$\Rightarrow$(v) is refuted by the unequal repeated-copy example described
below. The source route (ii)$\Rightarrow$(v) remains unproved after the
ZCL-derived common-weight reduction.

\paragraph{Amendment for the fixed-tensor implication in Theorem~4.9.}
The first convention is scale invariant at the level of the generated state
family, but it cannot be substituted unchanged into the implication to the
exact trace-preserving maps of Definition~4.1 in
Ref.~\cite{Cirac2016MPDO_arXiv}. The four-sector audit in
\path{docs/paper-gaps/cpgsv17_pf_rank_one.tex} exhibits the distinction. Its
raw tensor $\mathcal K$ satisfies
\begin{align}
    \mathcal T_{\mathcal K}^2
    &=\mathcal T_{\mathcal K},
    \label{eq:zcl_four_sector_raw_idempotence}
\end{align}
and has explicit renormalization maps, but its singleton BNT coefficient is
strictly subunit. The globally unit-weight representative
\begin{align}
    A
    &=\frac4{\sqrt5}\mathcal K
    \label{eq:zcl_four_sector_unit_representative}
\end{align}
satisfies \leanid{MPOTensor.IsSourceZCL} only up to scalar:
\begin{align}
    \mathcal T_A^2
    &=\frac4{\sqrt5}\mathcal T_A
    \label{eq:zcl_four_sector_scaled_relation}\\
    &\ne\mathcal T_A.
    \label{eq:zcl_four_sector_not_literal}
\end{align}
Its two-site blocking consequently fails Definition~4.1 of
Ref.~\cite{Cirac2016MPDO_arXiv}. The scalar absorption at source line~1660 is
performed only after literal ZCL has been assumed and does not transport
literal ZCL or the renormalization maps across rescaling. Thus the first
convention is the scale-invariant interpretation used by TNLean, whereas the
paper's fixed representative retains a separate literal diagrammatic
condition. Neither four-sector representative settles the exact
(ii)$\Rightarrow$(v) route.

A separate scalar BNT presentation with raw copy weights $1$ and $1/2$
satisfies the global normalization and every clause of condition~(iv), while
its two-site blocking fails Definition~4.1. The kernel-checked counterexample
is
\leanid{MPOTensor.CaseIIAbsorptionCounterexample.%
printed_theorem49_iv_to_v_is_false};
it refutes the literal implication (iv)$\Rightarrow$(v).

\section{The scale-invariant repaired interpretation and its consequences}

The repaired interpretation used in this development is the first convention.
In terms of the physical-trace transfer from
Eq.~(\ref{eq:zcl_physical_trace_transfer}), the predicate
\leanid{MPOTensor.IsSourceZCL} requires
\begin{align}
    \mathcal T_M
    &\ne0,
    \label{eq:zcl_scale_invariant_nonzero}\\
    \mathcal T_M^2
    &=\lambda\mathcal T_M
    &&\text{for some $\lambda\in\mathbb R$ with $\lambda>0$}.
    \label{eq:zcl_scale_invariant_quasi_idempotence}
\end{align}
The nonzero clause excludes the degenerate zero tensor, which would otherwise
satisfy the scalar relation vacuously for every $\lambda$. No spectral-radius
apparatus is required. If
$\mathcal T_M^2=\lambda\mathcal T_M$ with $\lambda>0$ and
$\mathcal T_M\ne0$, then the minimal polynomial of $\mathcal T_M$ divides
\begin{align}
    x(x-\lambda).
    \label{eq:zcl_minimal_polynomial_divisor}
\end{align}
Its eigenvalues therefore lie in
\begin{align}
    \{0,\lambda\},
    \label{eq:zcl_quasi_idempotent_spectrum}
\end{align}
and $\lambda$ is the leading eigenvalue. Literal idempotence is exactly the
case $\lambda=1$, corresponding to the physical-trace-normalized
representative. The maximally mixed witness has $\mathcal T_M=1$ and is
therefore classified as ZCL by the source transfer, whereas the doubled-index
condition excludes it.

\paragraph{The correction is to ZCL, not to the doubled-index ZCL predicate.}
The doubled-index transfer map $\mathcal E_M$ remains the object used by
\leanid{MPOTensor.IsZCL} and by the transfer-map formulation of the fusion and
algebra statements. It must not be identified with the source ZCL diagram.
Literal physical-trace idempotence implies normalized ZCL only together with
$\mathcal T_M\ne0$, or an equivalent generated-MPDO nondegeneracy assumption.
Without that clause, the zero transfer is idempotent but does not satisfy
Eqs.~(\ref{eq:zcl_scale_invariant_nonzero}) and
(\ref{eq:zcl_scale_invariant_quasi_idempotence}). Every source-facing
description that calls the doubled-index literal notion ``definitionally zero
correlation length'' must be weakened to a separate bridge with the nonzero
hypothesis stated explicitly.

\paragraph{Definition~4.1 gives the forward physical-trace equation.}
The theorem
\leanid{MPOTensor.isPhysicalTraceIdempotent_of_isRFPViaTS}, using the
physical-trace-square calculation, formalizes the ZCL part of
Proposition~\texttt{propsimple} in Appendix~C of
Ref.~\cite{Cirac2016MPDO_arXiv}, at source lines~1333--1340. That proposition
proves the first implication of Theorem~4.9 by establishing both ZCL and SAL.
At the explicit density-family boundary,
\leanid{MPOTensor.isSourceZCL_and_isSAL_of_isRFPViaTS_of_trace_ne_zero}
combines the physical-trace theorem with positive-length nonvanishing and
\leanid{MPOTensor.isSAL_of_isRFPViaTS_of_trace_ne_zero}. The horizontal
declarations are corollaries obtained by deriving that nonvanishing.

For every virtual matrix $X$, trace preservation of the one-to-two map gives
\begin{align}
    \tr(\mathcal T_M^2X)
    &=\tr(M_2(X))
    \label{eq:zcl_trace_pairing_two_site}\\
    &=\tr(M_1(X))
    \label{eq:zcl_trace_preservation_step}\\
    &=\tr(\mathcal T_MX).
    \label{eq:zcl_trace_pairing_one_site}
\end{align}
Nondegeneracy of the trace pairing yields
Eq.~(\ref{eq:zcl_literal_physical_idempotence}). Together with
$\mathcal T_M\ne0$, this implies
\leanid{MPOTensor.IsSourceZCL}. The explicit nonzero hypothesis is necessary
because the bare predicate \leanid{IsRFPViaTS} omits the source's standing
canonical-form and normalized-MPDO assumptions.

\paragraph{Effect on the purification witness.}
The diagonal purification in Eq.~(\ref{eq:zcl_purification_tensor}) has
$\mathcal T_M=1$ and therefore already satisfies the source ZCL equation. The
earlier claim that this witness fails zero correlation length resulted from
using $\mathcal E_M=\tfrac12\operatorname{id}$ in place of the physical-trace
transfer. Under the scale-invariant repaired interpretation, it is a positive
example showing that the up-to-scalar condition captures the maximally mixed
product state. This conclusion does not identify that condition with the
paper's literal fixed-representative diagram.

\paragraph{Historical resolution of the doubled-index coercion.}
Commit \texttt{43fb8749} introduced the former blocked-RFP construction, whose
proof treated doubled-index idempotence as the zero-correlation-length clause
of Theorem~4.9 in Ref.~\cite{Cirac2016MPDO_arXiv}. Commit
\texttt{74fb83dd} marked that dependence as unfaithful. Commit
\texttt{32482c47} removed the coercion and the paper-facing theorem chain, and
commit \texttt{fc135934} removed the remaining obsolete module. No such
coercion route remains in the development.

The present formalization separates the three transfer notions. The literal
Definition~4.2 conclusion
\begin{align}
    \mathcal T_M^2
    &=\mathcal T_M
    \label{eq:zcl_definition42_literal_conclusion}
\end{align}
follows directly from Definition~4.1 in
\leanid{MPOTensor.physTraceTransfer_sq_of_isRFPViaTS}. The broader
scale-invariant relation is recorded by \leanid{MPOTensor.IsSourceZCL}, while
\leanid{MPOTensor.IsZCL} concerns a different doubled-index transfer and is
not used as the source's ZCL clause. The doubled-index results in
\path{FusionIsometries.lean}, \path{AlgebraStructure.lean}, and
\path{CommutingFormBridge.lean} survive only as separately scoped conditional
helpers. None is a relocated proof of Theorem~4.9.

The normal Case-I structural conclusion is proved by
\leanid{MPOTensor.exists_physicalSectorFactorization_rank_one_coefficients_of_isSAL_of_literal_ZCL}.
Given one such factorization,
\leanid{MPOTensor.PhysicalSectorFactorization.NeighboringTraceFactorization.blockTwo_isRFPViaTS}
composes the two physical channels of Proposition~C.7 in
Ref.~\cite{Cirac2016MPDO_arXiv} twice, as observed at source
lines~1821--1825, and transports them back to the original physical
coordinates.

The remaining obligations belong to the stronger route from condition~(ii) of
Theorem~4.9 in Ref.~\cite{Cirac2016MPDO_arXiv}. First, the Case-II argument
needs a replacement for the false claim that coefficient absorption preserves
normality. This is the \emph{missing statement} in the
(ii)$\Rightarrow$(iv) node and is analyzed in
\path{docs/paper-gaps/cpgsv17_pf_rank_one.tex}.

Second, the Case-II argument at lines~1646--1665 of
Ref.~\cite{Cirac2016MPDO_arXiv} uses literal ZCL to make repeated-copy weights
independent of the copy label and absorb their common value. Lemma~C.10 is
linked exactly by
\leanid{MPOTensor.weight_copy_independent_of_isPhysicalTraceIdempotent}, with
its simple-canonical existential form supplied by
\leanid{MPOTensor.IsSimpleCanonicalForm.exists_weight_copy_independent_of_isPhysicalTraceIdempotent}.
Both results follow directly from the literal weighted-block equation. The
older \leanid{IsSourceZCL} declarations remain separate scale-invariant
generalizations and are no longer used as source nodes. After this
normalization, the Proposition~C.7 channel pairs must still be combined,
using the mutually orthogonal outer physical supports, into a single pair of
trace-preserving completely positive maps. This is the \emph{missing
statement} in the (ii)$\Rightarrow$(v) node.

By contrast, the raw implication (iv)$\Rightarrow$(v) is not a remaining
assembly problem. It is a refuted \emph{unfaithful statement}, and treating
the one-factorization result as though outer-sector assembly would prove it is
\emph{proof-path drift}. The hidden use of ZCL through
\texttt{lemmus} in the proof of \texttt{prop3to4} in
Ref.~\cite{Cirac2016MPDO_arXiv} is a separate instance of the same proof-path
drift. The former doubled-index theorem was removed rather than renamed. None
of the surviving gaps is a doubled-index bridge or a change of terminology.

\bibliographystyle{alpha}
\bibliography{references}

\end{document}
